\scanpage{65}
% archive foliation: 64; reused typescript, upside down in scan.

\begin{flushright}6.\end{flushright}

3.3. \uline{Les objectifs de recherche spécifiques des établissements d'enseignement supérieur.}

3.3.1. \uline{Mémorisation et diffusion de la connaissance} :

- rassemblement de la connaissance ; conservation de la connaissance ; mise en forme et diffusion de la connaissance. Cette diffusion ne sera pas limitée aux seuls étudiants à plein temps mais étendue vers toute personne intéressée et produite par tous supports(articles, ouvrages, collections, expositions) et toutes prestations (aide technique à l'industrie, études ...). La recherche pédagogique entre dans ce cadre.

3.3.2. \uline{Contribution à l'accroissement de la connaissance}

- en collaborant avec d'autres établissements de recherche, dans le cadre d'associations : l'établissement d'enseignement supérieur apporte une part des moyens, et une part importante du personnel compétent.

- en maintenant un état de ``fermentation intellectuelle'' permettant de faire naître des nouveaux domaines de recherche, notamment par le contact entre les diverses disciplines\unsure{.|;} Les nouvelles idées de recherche\inserted{,} initiées par l'établissement\inserted{,} doivent\inserted{,} ensuite être développées dans le cadre d'associations.

3.3.3. \uline{Formation par la recherche des futurs cadres de la nation.}

- Ceci exige un couplage de l'établissement avec des formations de recherche, où les étudiants en fin d'études pourront aller s'initier à la recherche en préparant une thèse de spécialité. La formation initiale et permanente des enseignants des établissements d'enseignement supérieur entre dans ce cadre.

3.3.4. \uline{Contribution à l'animation intellectuelle et socio-économique de la région} :

- Les universités sont en effet les seuls établissements à vocation de recherche qui n'ont pas leur siège à Paris, et dont le programme n'est pas commandé par l'autorité centrale. Elles doivent plus que jamais s'ouvrir aux préoccupations de leur région en prenant en

